% =====================================================================
% PEIB (SEES0082, UCL SSEES): equation slides for the tutorial handouts
% One \section per week. Frames are [plain], labelled wN_<topic>, in \Large,
% so each page exports as one cropped image for the Word handout.
% Frames follow the order of the handout, and the equation tags run in
% that order, so the images can be pasted top to bottom.
%
% Build (skill latex-teaching, see ~/.claude/skills/latex-teaching/SKILL.md):
%   py "$HOME/.claude/skills/latex-teaching/scripts/build.py" peib_main.tex --runs 1
%   or compile on Overleaf (upload this file alone) and save peib_main.pdf here.
% Export one PNG per page (0.83 = \Large shown as 12 pt in Word), one line.
% Student version (\showanswers = 0, 13 pages):
%   py "$HOME/.claude/skills/latex-teaching/scripts/export_frames.py" peib_main.pdf --out png/week2 --size-factor 0.83 --names w2_wt1_firm,w2_wt1_industry,w2_wt2_innovation,w2_wt3_equilibrium,w2_wt3_welfare,w2_wt4_elasticity,w2_wt5_firm,w2_wt5_industry,w2_wt6_derived,w2_wt7_chain,w2_appA_forms,w2_appC_passthrough,w2_appC_derived
% Tutor version (\showanswers = 1, 34 pages, tutor frames interleaved
% after the question they answer; file names = labels in file order):
%   py "$HOME/.claude/skills/latex-teaching/scripts/export_frames.py" peib_main.pdf --out png/week2_tutor --size-factor 0.83 --names w2_ans_q11,w2_ans_q12a,w2_ans_q12b,w2_wt1_firm,w2_wt1_industry,w2_wt2_innovation,w2_wt3_equilibrium,w2_wt3_welfare,w2_ans_q21a,w2_ans_q21b,w2_ans_q22a,w2_ans_q22c,w2_ans_q22b,w2_wt4_elasticity,w2_ans_q23a,w2_ans_q23c,w2_wt5_firm,w2_wt5_industry,w2_ans_q31b,w2_wt6_derived,w2_ans_q32a,w2_ans_q32b,w2_ans_q32n1,w2_ans_q32n2,w2_wt7_chain,w2_ans_q33a0,w2_ans_q33a1,w2_ans_q33b,w2_ans_q33c,w2_ans_q33c1,w2_ans_q33c2,w2_appA_forms,w2_appC_passthrough,w2_appC_derived
% Week 3 (26 frames, no answer frames: model responses stay text):
%   py "$HOME/.claude/skills/latex-teaching/scripts/build.py" peib_main.tex --runs 1 --frames w3_setup_tc,w3_setup_demand,w3_wt1_mc,w3_wt1_rule,w3_wt2_equilibrium,w3_wt2_cs_def,w3_wt2_ps_def,w3_wt2_welfare,w3_wt2_ts,w3_cs_step1,w3_cs_step2,w3_cs_step3,w3_cs_step4,w3_cs_step5,w3_cs_step6,w3_learning_curve,w3_production,w3_production_log,w3_production_spill,w3_cost_E,w3_wt3_experience,w3_wt3_cost,w3_wt4_supply,w3_wt5_ledger,w3_appA_incidence,w3_appA_example
%   py "$HOME/.claude/skills/latex-teaching/scripts/export_frames.py" peib_main_frames.pdf --out png/week3 --size-factor 0.83 --names w3_setup_tc,w3_setup_demand,w3_wt1_mc,w3_wt1_rule,w3_wt2_equilibrium,w3_wt2_cs_def,w3_wt2_ps_def,w3_wt2_welfare,w3_wt2_ts,w3_cs_step1,w3_cs_step2,w3_cs_step3,w3_cs_step4,w3_cs_step5,w3_cs_step6,w3_learning_curve,w3_production,w3_production_log,w3_production_spill,w3_cost_E,w3_wt3_experience,w3_wt3_cost,w3_wt4_supply,w3_wt5_ledger,w3_appA_incidence,w3_appA_example
% Compile one week only: uncomment the matching \includeonlyframes line below.
%
% Student vs tutor version: frames inside \ifnum\showanswers=1 (labels
% w2_ans_<question>, tags (M1), (M2), ...) reproduce the algebra of the
% model responses; keep \showanswers at 0 when exporting images for the
% student handout. Student tags (1), (2), ... are identical in both
% versions. Student frames only contain formulas that appear in the
% student-facing text (walkthroughs, appendices); anything from a model
% response is a tutor frame.
%
% Source of the week 2 content: handout version 4 (2026-10-12, 22:54),
% seminar2_handout.docx (student master) and
% seminar2_handout_with_model_responses.docx (generated from it by
% make_tutor.py). Paragraph indices [n] in the frame comments refer to
% the student text export (student3_text.md), [Tn] to the tutor export
% (tutor3_text.md). The handout keeps 13 equation image slots, one per
% student frame, in this order (the image follows the paragraph named):
%   w2_wt1_firm          Walkthrough 1, Step 1 ("... so we solve the same equation for q")
%   w2_wt1_industry      Walkthrough 1, Step 3 ("... each moving only a little up its MC curve")
%   w2_wt2_innovation    Walkthrough 2 ("Repeating Steps 1 to 3 with the new cost function ...")
%   w2_wt3_equilibrium   Walkthrough 3 ("The market clears where ...")
%   w2_wt3_welfare       Walkthrough 3 ("With linear curves, consumer and producer surplus ...")
%   w2_wt4_elasticity    Walkthrough 4 ("Elasticity: how responsive quantity is to price ...")
%   w2_wt5_firm          Walkthrough 5 ("Read only. The algebra is the same as in Question 1.1 ...")
%   w2_wt5_industry      Walkthrough 5 (directly after w2_wt5_firm)
%   w2_wt6_derived       Walkthrough 6 ("For any nickel price PN, the battery market clears ...")
%   w2_wt7_chain         Walkthrough 7 ("The innovation travels down the chain in two stages ...")
%   w2_appA_forms        Appendix A ("Moving between the two forms is simple algebra ...")
%   w2_appC_passthrough  Appendix C ("Pass-through in elasticities ...")
%   w2_appC_derived      Appendix C ("Repeat the derivation of Walkthrough 6 in general terms ...")
% The tutor handout has one more slot, inside the model response to
% Question 2.3 between parts (b) and (c): frame w2_ans_q23c.
% Retired frames (refresher, old worked examples, review-note frames) are
% in archive.tex next to this file.
% Notation follows the handout: Walkthrough 4 writes P = A - bQ and
% P = u + sQ; from Walkthrough 7 on, b is the slope of battery demand
% (= 1) and g the slope of other battery marginal cost (= 1), so the
% nickel demand slope is b + g = 2.
% =====================================================================

\documentclass{beamer}

\usepackage{tikz}
\usepackage{pgfplots}
\usepackage{amsmath}
\usepackage{amssymb}
\usepackage{xcolor}
\usepackage{booktabs}
\usepackage{colortbl}
\usepackage{array}

\usetikzlibrary{decorations.pathreplacing}
\pgfplotsset{compat=newest}

% Minimal, no headers/footers/nav symbols
\mode<presentation>{
  \usetheme{default}
  \setbeamertemplate{navigation symbols}{}
  \setbeamertemplate{headline}{}
  \setbeamertemplate{footline}{}
}

% Custom colors (international deck)
\definecolor{myblue}{RGB}{39,64,139}
\definecolor{homecolor}{RGB}{0,102,204}
\definecolor{foreigncolor}{RGB}{204,51,0}
\definecolor{capitalcolor}{RGB}{0,102,204}
\definecolor{laborcolor}{RGB}{204,51,0}
\definecolor{landcolor}{RGB}{34,139,34}

% PEIB website palette (peib_cols in Teaching_PEIB/code/functions_models.R)
\definecolor{peibdemand}{HTML}{1F3B5C}
\definecolor{peibsupply}{HTML}{8C4A00}
\definecolor{peibcs}{HTML}{4C78A8}
\definecolor{peibps}{HTML}{F58518}
\definecolor{peibtax}{HTML}{54A24B}
\definecolor{peibdwl}{HTML}{E45756}
\definecolor{peibtr}{HTML}{72B7B2}
\definecolor{peibtvc}{HTML}{B279A2}
\definecolor{peibprofit}{HTML}{EECA3B}
\definecolor{peibfc}{HTML}{BAB0AC}

% ANSWERS TOGGLE: 1 = include tutor-only frames, 0 = student version
\newcommand{\showanswers}{0}

% SELECTIVE COMPILATION: uncomment one line to build a single week
% \includeonlyframes{w2_ans_q11,w2_ans_q12a,w2_ans_q12b,w2_wt1_firm,w2_wt1_industry,w2_wt2_innovation,w2_wt3_equilibrium,w2_wt3_welfare,w2_ans_q21a,w2_ans_q21b,w2_ans_q22a,w2_ans_q22c,w2_ans_q22b,w2_wt4_elasticity,w2_ans_q23a,w2_ans_q23c,w2_wt5_firm,w2_wt5_industry,w2_ans_q31b,w2_wt6_derived,w2_ans_q32a,w2_ans_q32b,w2_ans_q32n1,w2_ans_q32n2,w2_wt7_chain,w2_ans_q33a0,w2_ans_q33a1,w2_ans_q33b,w2_ans_q33c,w2_ans_q33c1,w2_ans_q33c2,w2_appA_forms,w2_appC_passthrough,w2_appC_derived}

\begin{document}

%========================
% WEEK 1: Competitive market, welfare (Tutorial 1)
%========================
\section{Week 1}

% No equation slides yet. Add frames here as [plain,label=w1_<topic>].


%========================
% WEEK 2: Linked markets, pass-through, welfare along a value chain
% (Tutorial 2). Handout structure: Exercise 1 (Questions 1.1-1.2,
% Walkthroughs 1-2), Exercise 2 (Walkthrough 3, Questions 2.1-2.2,
% Walkthrough 4, Question 2.3), Exercise 3 (Walkthrough 5, Question 3.1,
% Walkthrough 6, Question 3.2, Walkthrough 7, Question 3.3),
% Appendices A (two forms), B (refresher, text only), C (optional).
%========================
\section{Week 2}

% --- Tutor: Question 1.1 model responses [T57]-[T59] ------------------------
% Hand check: FC = 20; TVC = 50q + 5q^2; MC = 50 + 2*5q = 50 + 10q;
% 90 = 50 + 10q => 10q = 40 => q* = 4
\ifnum\showanswers=1
\begin{frame}[plain,label=w2_ans_q11]
  \Large
  \[
    \tag{M1}\;FC = 20, \qquad TVC(q) = 50q + 5q^2
  \]
  \[
    \tag{M2}\;MC(q) = \frac{dTC_N}{dq} = 50 + 10q
  \]
  \begin{align*}
    \tag{M3}\;90 &= 50 + 10q\\
    \Rightarrow 10q &= 40\\
    \Rightarrow q^* &= 4
  \end{align*}

  \vspace{0.3em}
  \begin{itemize}\Large
    \item $5q^2$ contributes $10q$ (not $5q$)
  \end{itemize}
\end{frame}
\fi

% --- Tutor: Question 1.2 a-b model responses [T63] --------------------------
% Hand check: TR = 90*4 = 360; TVC(4) = 200 + 5*16 = 200 + 80 = 280;
% trapezium 0.5*(50 + 90)*4 = 0.5*140*4 = 280
\ifnum\showanswers=1
\begin{frame}[plain,label=w2_ans_q12a]
  \Large
  \[
    \tag{M4}\;TR = P \times q^* = 90 \times 4 = 360
  \]
  \begin{align*}
    \tag{M5}\;TVC(4) &= 50 \times 4 + 5 \times 4^2\\
                     &= 200 + 80 = 280
  \end{align*}
  \[
    \tag{M6}\;\text{Check: } \tfrac{1}{2} \times (50 + 90) \times 4 = 280 \;\checkmark
  \]

  \vspace{0.3em}
  \begin{itemize}\Large
    \item $TVC$: area under $MC$ from $0$ to $4$ (a trapezium)
  \end{itemize}
\end{frame}
\fi

% --- Tutor: Question 1.2 c-e model responses [T63]-[T64] --------------------
% Hand check: PS = 360 - 280 = 80; triangle 0.5*4*40 = 80;
% profit = 360 - 280 - 20 = 60; PS = profit + FC: 80 = 60 + 20
\ifnum\showanswers=1
\begin{frame}[plain,label=w2_ans_q12b]
  \Large
  \[
    \tag{M7}\;PS = TR - TVC = 360 - 280 = 80
  \]
  \[
    \tag{M8}\;\text{Check: } \tfrac{1}{2} \times 4 \times (90 - 50) = 80 \;\checkmark
  \]
  \begin{align*}
    \tag{M9}\;\text{Profit} &= TR - TVC - FC\\
                            &= 360 - 280 - 20 = 60
  \end{align*}
  \[
    \tag{M10}\;PS = \text{Profit} + FC: \quad 80 = 60 + 20
  \]

  \vspace{0.3em}
  \begin{itemize}\Large
    \item $FC$ does not depend on $q$: not in the figure
  \end{itemize}
\end{frame}
\fi

% --- Walkthrough 1, Step 1 [56]: the supply function of one firm -----------
% Hand check: P = 50 + 10q => 10q = P - 50 => q = (P - 50)/10; q(50) = 0
\begin{frame}[plain,label=w2_wt1_firm]
  \Large
  \begin{align*}
    \tag{1}\;P &= MC(q) = 50 + 10q\\
    \Rightarrow 10q &= P - 50\\
    \Rightarrow q(P) &= \frac{P - 50}{10}, \quad P \ge 50
  \end{align*}

  \vspace{0.5em}
  \begin{itemize}\Large
    \item Firm's supply curve $=$ its $MC$ curve
    \item Direct form: $q$ as a function of $P$
    \item $P < 50$: no supply (first unit costs $50$)
  \end{itemize}
\end{frame}

% --- Walkthrough 1, Steps 2-3 [58]-[59]: the industry supply curve ---------
% Hand check: Q = 10*(P - 50)/10 = P - 50 => P = 50 + Q; firm slope 10,
% industry slope 1 (+1 pound: +0.1 per firm, +1 for 10 firms)
\begin{frame}[plain,label=w2_wt1_industry]
  \Large
  \begin{align*}
    \tag{2}\;Q_N(P) &= 10 \times q(P) = 10 \times \frac{P - 50}{10}\\
                    &= P - 50
  \end{align*}
  \[
    \tag{3}\;P_N = 50 + Q_N
  \]

  \vspace{0.5em}
  \begin{itemize}\Large
    \item Horizontal sum: add quantities, not prices
    \item Slope: firm $10$, industry $1$
    \item A \pounds 1 rise: $+0.1$ per firm, $+1$ in total
  \end{itemize}
\end{frame}

% --- Walkthrough 2 [65]-[68]: the innovation shifts the supply curves ------
% Hand check: MC = 20 + 10q; q = (P - 20)/10 (P >= 20); Q = P - 20,
% P = 20 + Q. Vertical shift 50 - 20 = 30 at every Q
\begin{frame}[plain,label=w2_wt2_innovation]
  \Large
  \[
    \tag{4}\;MC(q) = \frac{dTC_N}{dq} = 20 + 10q
  \]
  \[
    \tag{5}\;q(P) = \frac{P - 20}{10}, \quad P \ge 20
  \]
  \[
    \tag{6}\;Q_N = P - 20 \;\Leftrightarrow\; P_N = 20 + Q_N
  \]

  \vspace{0.5em}
  \begin{itemize}\Large
    \item Steps 1 to 3 with the new cost function
    \item Down by $30$ at every $Q$ (the cost change)
    \item A shift of supply, not a movement along it
  \end{itemize}
\end{frame}

% --- Walkthrough 3 [80]: equilibrium before the innovation -----------------
% Hand check: 170 - 2Q = 50 + Q => 120 = 3Q => Q = 40; P = 90;
% demand 170 - 80 = 90; 40/10 = 4 per firm
\begin{frame}[plain,label=w2_wt3_equilibrium]
  \Large
  \begin{align*}
    \tag{7}\;170 - 2Q &= 50 + Q\\
    \Rightarrow 3Q &= 120\\
    \Rightarrow Q_0 &= 40
  \end{align*}
  \begin{align*}
    \tag{8}\;P_0 &= 50 + 40 = 90\\
    \text{Check: } 170 - 2 \times 40 &= 90 \;\checkmark
  \end{align*}

  \vspace{0.3em}
  \begin{itemize}\Large
    \item Inverse form: set the two prices equal
    \item $E_0 = (40, 90)$; $4$ units per firm
  \end{itemize}
\end{frame}

% --- Walkthrough 3 [82]: welfare before the innovation ---------------------
% Hand check: CS = 0.5*40*80 = 1600; PS = 0.5*40*40 = 800; TS = 2400
% (the per-firm check 10 * 80 = 800 is tutor-only: PS = 80 answers Q1.2c)
\begin{frame}[plain,label=w2_wt3_welfare]
  \Large
  \begin{align*}
    \tag{9}\;CS_{N,0} &= \tfrac{1}{2} \times 40 \times (170 - 90)\\
                     &= \tfrac{1}{2} \times 40 \times 80 = 1{,}600
  \end{align*}
  \begin{align*}
    \tag{10}\;PS_{N,0} &= \tfrac{1}{2} \times 40 \times (90 - 50)\\
                      &= \tfrac{1}{2} \times 40 \times 40 = 800
  \end{align*}
  \[
    \tag{11}\;TS_{N,0} = 1{,}600 + 800 = 2{,}400
  \]

  \vspace{0.3em}
  \begin{itemize}\Large
    \item The two triangles in Figure 4
  \end{itemize}
\end{frame}

% --- Tutor: Question 2.1a model response [T102] -----------------------------
% Hand check: 170 - 2Q = 20 + Q => 150 = 3Q => Q = 50; P = 70;
% 170 - 100 = 70; q = (70 - 20)/10 = 5
\ifnum\showanswers=1
\begin{frame}[plain,label=w2_ans_q21a]
  \Large
  \begin{align*}
    \tag{M11}\;170 - 2Q &= 20 + Q\\
    \Rightarrow 3Q &= 150\\
    \Rightarrow Q_1 &= 50
  \end{align*}
  \begin{align*}
    \tag{M12}\;P_1 &= 20 + 50 = 70\\
    \text{Check: } 170 - 2 \times 50 &= 70 \;\checkmark
  \end{align*}
  \[
    \tag{M13}\;q = \frac{70 - 20}{10} = 5 \quad (\text{was } 4)
  \]
\end{frame}
\fi

% --- Tutor: Question 2.1a (direct form) and 2.1b [T102]-[T103] --------------
% Hand check: 85 - 0.5P = P - 20 => 105 = 1.5P => P = 70; Q = 50;
% dP = 70 - 90 = -20; dQ = 50 - 40 = 10; MC of 50th unit 20 + 50 = 70
\ifnum\showanswers=1
\begin{frame}[plain,label=w2_ans_q21b]
  \Large
  \begin{align*}
    \tag{M14}\;85 - 0.5P &= P - 20\\
    \Rightarrow 1.5P &= 105\\
    \Rightarrow P_1 &= 70
  \end{align*}
  \[
    \tag{M15}\;Q_1 = 70 - 20 = 50
  \]
  \begin{align*}
    \tag{M16}\;\Delta P &= 70 - 90 = -20\\
    \tag{M17}\;\Delta Q &= 50 - 40 = +10
  \end{align*}
  \begin{itemize}\Large
    \item $MC$ of the 50th unit: $20 + 50 = 70 = P_1$
  \end{itemize}
\end{frame}
\fi

% --- Tutor: Question 2.2a model response [T107] -----------------------------
% Hand check: CS = 0.5*50*(170 - 70) = 0.5*50*100 = 2500;
% PS = 0.5*50*(70 - 20) = 0.5*50*50 = 1250; TS = 3750
\ifnum\showanswers=1
\begin{frame}[plain,label=w2_ans_q22a]
  \Large
  \begin{align*}
    \tag{M18}\;CS_{N,1} &= \tfrac{1}{2} \times 50 \times (170 - 70)\\
                       &= \tfrac{1}{2} \times 50 \times 100 = 2{,}500
  \end{align*}
  \begin{align*}
    \tag{M19}\;PS_{N,1} &= \tfrac{1}{2} \times 50 \times (70 - 20)\\
                       &= \tfrac{1}{2} \times 50 \times 50 = 1{,}250
  \end{align*}
  \[
    \tag{M20}\;TS_{N,1} = 2{,}500 + 1{,}250 = 3{,}750
  \]
\end{frame}
\fi

% --- Tutor: Question 2.2a per-firm check [T107] -----------------------------
% Hand check: TR = 70*5 = 350; TVC = 20*5 + 5*25 = 100 + 125 = 225;
% PS = 125; 10*125 = 1250
\ifnum\showanswers=1
\begin{frame}[plain,label=w2_ans_q22c]
  \Large
  \[
    \tag{M21}\;TR = 70 \times 5 = 350
  \]
  \[
    \tag{M22}\;TVC = 20 \times 5 + 5 \times 5^2 = 225
  \]
  \[
    \tag{M23}\;PS = 350 - 225 = 125
  \]
  \[
    \tag{M24}\;10 \times 125 = 1{,}250 = PS_{N,1} \;\checkmark
  \]

  \vspace{0.3em}
  \begin{itemize}\Large
    \item Per firm, after the innovation ($q = 5$)
  \end{itemize}
\end{frame}
\fi

% --- Tutor: Question 2.2b-c model responses [T108]-[T109] -------------------
% Hand check: 2500 - 1600 = 900; 1250 - 800 = 450; 3750 - 2400 = 1350
\ifnum\showanswers=1
\begin{frame}[plain,label=w2_ans_q22b]
  \Large
  \begin{align*}
    \tag{M25}\;\Delta CS_N &= 2{,}500 - 1{,}600 = +900\\
    \tag{M26}\;\Delta PS_N &= 1{,}250 - 800 = +450\\
    \tag{M27}\;\Delta TS_N &= 3{,}750 - 2{,}400 = +1{,}350
  \end{align*}

  \vspace{0.3em}
  \begin{itemize}\Large
    \item Producers gain: cost $-30$, price only $-20$
    \item Nickel buyers: battery firms, not households
  \end{itemize}
\end{frame}
\fi

% --- Walkthrough 4 [88]: definition of the price elasticity ----------------
% Hand check: definition only; (dQ/dP)(P/Q) is the limit of (dQ/Q)/(dP/P)
\begin{frame}[plain,label=w2_wt4_elasticity]
  \Large
  \begin{align*}
    \tag{12}\;\varepsilon_D &= \frac{\%\Delta Q}{\%\Delta P}
      = \frac{\Delta Q / Q}{\Delta P / P}\\
    &= \frac{\Delta Q}{\Delta P} \times \frac{P}{Q}
  \end{align*}
  \[
    \tag{13}\;\varepsilon_D = \frac{dQ}{dP} \times \frac{P}{Q}
    \qquad \text{(point elasticity)}
  \]

  \vspace{0.3em}
  \begin{itemize}\Large
    \item $dQ/dP$: slope of the direct form $Q(P)$
    \item $P/Q$ cancels the units: a pure number
    \item $\varepsilon_S$: same formula, positive
    \item $|\varepsilon| > 1$ elastic, $|\varepsilon| < 1$ inelastic
  \end{itemize}
\end{frame}

% --- Tutor: Question 2.3a model response [T124] -----------------------------
% Hand check: 70/50 = 1.4; eps_D = -0.5*1.4 = -0.7; eps_S = 1*1.4 = 1.4;
% P/Q at E0 = 90/40 = 2.25
\ifnum\showanswers=1
\begin{frame}[plain,label=w2_ans_q23a]
  \Large
  \[
    \tag{M28}\;\varepsilon_D = -0.5 \times \frac{70}{50} = -0.7
  \]
  \[
    \tag{M29}\;\varepsilon_S = 1 \times \frac{70}{50} = 1.4
  \]
  \[
    \tag{M30}\;\frac{P}{Q}: \quad \frac{90}{40} = 2.25 \;\longrightarrow\; \frac{70}{50} = 1.4
  \]

  \vspace{0.3em}
  \begin{itemize}\Large
    \item Slopes unchanged; $P/Q$ has fallen
    \item Demand now inelastic: $0.7 < 1$
  \end{itemize}
\end{frame}
\fi

% --- Tutor: Question 2.3c model response [T127]; image slot [T126] ---------
% Hand check: b/(b + s) = 1/(1 + s/b) -> 1 as b -> infinity; b = 0: 0/s = 0;
% s = 0: b/b = 1; b = 0: Q = 90 - 20 = 70 (+30 units)
\ifnum\showanswers=1
\begin{frame}[plain,label=w2_ans_q23c]
  \Large
  \[
    \tag{M31}\;\rho = \frac{b}{b + s} = \frac{1}{1 + s/b} \xrightarrow{\;b \to \infty\;} 1
  \]
  \[
    \tag{M32}\;b = 0: \quad \rho = \frac{0}{0 + s} = 0
  \]
  \[
    \tag{M33}\;s = 0: \quad \rho = \frac{b}{b + 0} = 1
  \]

  \vspace{0.3em}
  \begin{itemize}\Large
    \item $b \to \infty$: vertical demand, $Q$ fixed
    \item $b = 0$: horizontal demand, $P$ fixed
    \item $s = 0$: flat supply, constant $MC$
  \end{itemize}
\end{frame}
\fi

% --- Walkthrough 5 [108]: the battery firm's marginal cost and supply ------
% Hand check: d/dq [20 + (P_N + 50)q + 5q^2] = P_N + 50 + 10q;
% P_B = P_N + 50 + 10q => 10q = P_B - P_N - 50 => q = (P_B - P_N - 50)/10
\begin{frame}[plain,label=w2_wt5_firm]
  \Large
  \[
    \tag{14}\;MC_B(q) = \frac{dTC_B}{dq} = P_N + 50 + 10q
  \]
  \begin{align*}
    \tag{15}\;P_B &= P_N + 50 + 10q\\
    \Rightarrow 10q &= P_B - P_N - 50\\
    \Rightarrow q &= \frac{P_B - P_N - 50}{10}
  \end{align*}

  \vspace{0.3em}
  \begin{itemize}\Large
    \item $P_N + 50$ replaces the $50$ of Exercise 1
    \item Valid for $P_B \ge P_N + 50$
  \end{itemize}
\end{frame}

% --- Walkthrough 5 [108], [111]: the battery industry supply curve ---------
% Hand check: Q = 10*(P_B - P_N - 50)/10 = P_B - P_N - 50 => P_B = P_N + 50 + Q
\begin{frame}[plain,label=w2_wt5_industry]
  \Large
  \begin{align*}
    \tag{16}\;Q_B &= 10 \times \frac{P_B - P_N - 50}{10}\\
                 &= P_B - P_N - 50
  \end{align*}
  \[
    \tag{17}\;P_B = P_N + 50 + Q_B
  \]

  \vspace{0.3em}
  \begin{itemize}\Large
    \item Intercept $P_N + 50$: nickel plus other cost
    \item Slope $1$: rising other marginal cost
    \item One curve for each nickel price
  \end{itemize}
\end{frame}

% --- Tutor: Question 3.1b model response [T146] -----------------------------
% Hand check: (P_N - 1) + 50 = (P_N + 50) - 1, intercept down by 1;
% P_N 90 -> 70: intercept 140 -> 120 (down by 20)
\ifnum\showanswers=1
\begin{frame}[plain,label=w2_ans_q31b]
  \Large
  \[
    \tag{M34}\;P_B = (P_N - 1) + 50 + Q_B
  \]
  \begin{align*}
    \tag{M35}\;P_N = 90: \quad P_B &= 140 + Q_B\\
    P_N = 70: \quad P_B &= 120 + Q_B
  \end{align*}

  \vspace{0.3em}
  \begin{itemize}\Large
    \item $P_N$ falls by \pounds 1: curve shifts down by \pounds 1
    \item $P_N$ from $90$ to $70$: down by $20$
    \item Demand unchanged: $P_B = 220 - Q_B$
  \end{itemize}
\end{frame}
\fi

% --- Walkthrough 6 [117]: deriving the demand for nickel -------------------
% Hand check: 220 - Q = P_N + 50 + Q => 170 - P_N = 2Q => P_N = 170 - 2Q;
% 2Q = 170 - P_N => Q = 85 - 0.5P_N
\begin{frame}[plain,label=w2_wt6_derived]
  \Large
  \begin{align*}
    \tag{18}\;220 - Q_B &= P_N + 50 + Q_B\\
    \Rightarrow 2Q_B &= 170 - P_N\\
    \Rightarrow P_N &= 170 - 2Q_B = 170 - 2Q_N
  \end{align*}
  \[
    \tag{19}\;\text{Direct form: } Q_N = 85 - 0.5P_N
  \]

  \vspace{0.3em}
  \begin{itemize}\Large
    \item The nickel demand curve of Exercise 2
    \item $Q_N = Q_B$: one unit of nickel per battery
  \end{itemize}
\end{frame}

% --- Tutor: Question 3.2a-b before [T161] -----------------------------------
% Hand check: 220 - 40 = 180; 90 + 50 + 40 = 180; q = (180 - 90 - 50)/10 = 40/10 = 4
\ifnum\showanswers=1
\begin{frame}[plain,label=w2_ans_q32a]
  \Large
  \begin{align*}
    \tag{M36}\;P_{B,0} &= 220 - 40 = 180\\
    \text{Check: } 90 + 50 + 40 &= 180 \;\checkmark
  \end{align*}
  \[
    \tag{M37}\;q_0 = \frac{180 - 90 - 50}{10} = 4
  \]

  \vspace{0.3em}
  \begin{itemize}\Large
    \item Before: $Q_B = 40$, $P_N = 90$
  \end{itemize}
\end{frame}
\fi

% --- Tutor: Question 3.2a-c after [T161] ------------------------------------
% Hand check: 220 - 50 = 170; 70 + 50 + 50 = 170; q = (170 - 70 - 50)/10 = 50/10 = 5
\ifnum\showanswers=1
\begin{frame}[plain,label=w2_ans_q32b]
  \Large
  \begin{align*}
    \tag{M38}\;P_{B,1} &= 220 - 50 = 170\\
    \text{Check: } 70 + 50 + 50 &= 170 \;\checkmark
  \end{align*}
  \[
    \tag{M39}\;q_1 = \frac{170 - 70 - 50}{10} = 5
  \]

  \vspace{0.3em}
  \begin{itemize}\Large
    \item After: $Q_B = 50$, $P_N = 70$
    \item Each battery firm: from $4$ to $5$ units
    \item $P_B$: $180 \to 170$, as in Walkthrough 7
  \end{itemize}
\end{frame}
\fi

% --- Tutor: Question 3.2, tutor note (general formulae) [T163] -------------
% Hand check: (220 - 50 - 50)/3 = 120/3 = 40 (u = 50); k/(b + g + s) = 30/3 = 10
\ifnum\showanswers=1
\begin{frame}[plain,label=w2_ans_q32n1]
  \Large
  \begin{align*}
    \tag{M40}\;Q &= \frac{A - c - u}{b + g + s}\\
                 &= \frac{220 - 50 - 50}{3} = 40
  \end{align*}
  \[
    \tag{M41}\;\Delta Q = \frac{k}{b + g + s} = \frac{30}{3} = 10
  \]

  \vspace{0.3em}
  \begin{itemize}\Large
    \item Nickel supply $P_N = u + sQ$, cost cut $k$
    \item $b = g = s = 1$, $k = 30$
  \end{itemize}
\end{frame}
\fi

% --- Tutor: Question 3.2, tutor note (general formulae) [T163] -------------
% Hand check: -30*2/3 = -20; -30*1/3 = -10; rates 2/3, 1/2, 1/3
\ifnum\showanswers=1
\begin{frame}[plain,label=w2_ans_q32n2]
  \Large
  \begin{align*}
    \tag{M42}\;\Delta P_N &= -k\,\frac{b + g}{b + g + s}\\
                          &= -30 \times \frac{2}{3} = -20
  \end{align*}
  \begin{align*}
    \tag{M43}\;\Delta P_B &= -k\,\frac{b}{b + g + s}\\
                          &= -30 \times \frac{1}{3} = -10
  \end{align*}

  \vspace{0.3em}
  \begin{itemize}\Large
    \item Rates: $(b+g)/(b+g+s)$, $b/(b+g)$, $b/(b+g+s)$
  \end{itemize}
\end{frame}
\fi

% --- Walkthrough 7 [127]: the pass-through chain ---------------------------
% Hand check: 20/30 = 2/3 = 2/(2 + 1); 10/20 = 1/2 = 1/(1 + 1);
% 10/30 = 1/3 = 2/3 * 1/2; P_B: 180 -> 170
\begin{frame}[plain,label=w2_wt7_chain]
  \Large
  \[
    \tag{20}\;\Delta MC_N = -30 \;\longrightarrow\; \Delta P_N = -20
  \]
  \[
    \tag{21}\;\Delta P_N = -20 \;\longrightarrow\; \Delta MC_B = -20
  \]
  \[
    \tag{22}\;\Delta MC_B = -20 \;\longrightarrow\; \Delta P_B = -10
  \]
  \[
    \tag{23}\;\rho = \frac{2}{3} \times \frac{1}{2} = \frac{10}{30} = \frac{1}{3}
  \]

  \vspace{0.3em}
  \begin{itemize}\Large
    \item Upstream $2/3$; downstream $b/(b + g) = 1/2$
    \item $\Delta MC_B = \Delta P_N$: one unit of nickel per battery
  \end{itemize}
\end{frame}

% --- Tutor: Question 3.3a before [T178] --------------------------------------
% Hand check: 0.5*40*(220 - 180) = 0.5*40*40 = 800;
% intercept 90 + 50 = 140; 0.5*40*(180 - 140) = 800
\ifnum\showanswers=1
\begin{frame}[plain,label=w2_ans_q33a0]
  \Large
  \begin{align*}
    \tag{M44}\;CS_{B,0} &= \tfrac{1}{2} \times 40 \times (220 - 180)\\
                       &= \tfrac{1}{2} \times 40 \times 40 = 800
  \end{align*}
  \begin{align*}
    \tag{M45}\;PS_{B,0} &= \tfrac{1}{2} \times 40 \times (180 - 140)\\
                       &= \tfrac{1}{2} \times 40 \times 40 = 800
  \end{align*}

  \vspace{0.3em}
  \begin{itemize}\Large
    \item Supply intercept $140 = 90 + 50$ (not $50$)
  \end{itemize}
\end{frame}
\fi

% --- Tutor: Question 3.3a after [T178] ---------------------------------------
% Hand check: 0.5*50*(220 - 170) = 0.5*50*50 = 1250;
% intercept 70 + 50 = 120; 0.5*50*(170 - 120) = 1250
\ifnum\showanswers=1
\begin{frame}[plain,label=w2_ans_q33a1]
  \Large
  \begin{align*}
    \tag{M46}\;CS_{B,1} &= \tfrac{1}{2} \times 50 \times (220 - 170)\\
                       &= \tfrac{1}{2} \times 50 \times 50 = 1{,}250
  \end{align*}
  \begin{align*}
    \tag{M47}\;PS_{B,1} &= \tfrac{1}{2} \times 50 \times (170 - 120)\\
                       &= \tfrac{1}{2} \times 50 \times 50 = 1{,}250
  \end{align*}

  \vspace{0.3em}
  \begin{itemize}\Large
    \item Supply intercept $120 = 70 + 50$
  \end{itemize}
\end{frame}
\fi

% --- Tutor: Question 3.3b and 3.3d: the ledger [T181], [T184] ---------------
% Hand check: 800*3 = 2400; 1250*3 = 3750; 3750 - 2400 = 1350; 1350/3 = 450
\ifnum\showanswers=1
\begin{frame}[plain,label=w2_ans_q33b]
  \Large
  \begin{align*}
    \tag{M48}\;NSB_0 &= PS_N + PS_B + CS_B\\
                    &= 800 + 800 + 800 = 2{,}400
  \end{align*}
  \[
    \tag{M49}\;NSB_1 = 3 \times 1{,}250 = 3{,}750
  \]
  \[
    \tag{M50}\;\Delta NSB = 3{,}750 - 2{,}400 = +1{,}350
  \]

  \vspace{0.3em}
  \begin{itemize}\Large
    \item Equals $TS_N$ of Exercise 2
    \item Split: $+450$ to each group
  \end{itemize}
\end{frame}
\fi

% --- Tutor: Question 3.3c: no double counting [T182] -------------------------
% Hand check: 800 + 800 = 1600; 1250 + 1250 = 2500; 1600 + 2400 = 4000
\ifnum\showanswers=1
\begin{frame}[plain,label=w2_ans_q33c]
  \Large
  \[
    \tag{M51}\;CS_N = CS_B + PS_B
  \]
  \begin{align*}
    \tag{M52}\;1{,}600 &= 800 + 800\\
    \tag{M53}\;2{,}500 &= 1{,}250 + 1{,}250
  \end{align*}
  \[
    \tag{M54}\;\text{Wrong: } CS_N + NSB_0 = 4{,}000
  \]

  \vspace{0.3em}
  \begin{itemize}\Large
    \item Count each benefit once
    \item $CS_N + PS_N$ or $PS_N + PS_B + CS_B$
  \end{itemize}
\end{frame}
\fi

% --- Tutor: Question 3.3c, check via resources, before [T183] ---------------
% Hand check: 220*40 - 0.5*1600 = 8800 - 800 = 8000;
% 50*40 + 800 = 2800 (twice); 8000 - 2800 - 2800 = 2400
\ifnum\showanswers=1
\begin{frame}[plain,label=w2_ans_q33c1]
  \Large
  \begin{align*}
    \tag{M55}\;V &= 220 \times 40 - \tfrac{1}{2} \times 40^2 = 8{,}000\\
    \tag{M56}\;C_N &= 50 \times 40 + \tfrac{1}{2} \times 40^2 = 2{,}800\\
    \tag{M57}\;C_O &= 50 \times 40 + \tfrac{1}{2} \times 40^2 = 2{,}800
  \end{align*}
  \[
    \tag{M58}\;8{,}000 - 2{,}800 - 2{,}800 = 2{,}400 \;\checkmark
  \]

  \vspace{0.3em}
  \begin{itemize}\Large
    \item $V$: value to final users
    \item $C_N$, $C_O$: nickel and other cost
    \item $P_N \times Q$: a transfer, not a cost
  \end{itemize}
\end{frame}
\fi

% --- Tutor: Question 3.3c, check via resources, after [T183] ----------------
% Hand check: 220*50 - 0.5*2500 = 11000 - 1250 = 9750;
% 20*50 + 1250 = 2250; 50*50 + 1250 = 3750; 9750 - 2250 - 3750 = 3750
\ifnum\showanswers=1
\begin{frame}[plain,label=w2_ans_q33c2]
  \Large
  \begin{align*}
    \tag{M59}\;V &= 220 \times 50 - \tfrac{1}{2} \times 50^2 = 9{,}750\\
    \tag{M60}\;C_N &= 20 \times 50 + \tfrac{1}{2} \times 50^2 = 2{,}250\\
    \tag{M61}\;C_O &= 50 \times 50 + \tfrac{1}{2} \times 50^2 = 3{,}750
  \end{align*}
  \[
    \tag{M62}\;9{,}750 - 2{,}250 - 3{,}750 = 3{,}750 \;\checkmark
  \]

  \vspace{0.3em}
  \begin{itemize}\Large
    \item After: nickel $MC = 20 + Q$, $Q = 50$
  \end{itemize}
\end{frame}
\fi

% --- Appendix A [143]: direct and inverse form -----------------------------
% Hand check: Q = P - 50 <=> P = 50 + Q; Q = 85 - 0.5P <=> 0.5P = 85 - Q
% <=> P = 170 - 2Q. Slopes: demand dQ/dP = -0.5, dP/dQ = -2 = 1/(-0.5);
% supply 1 and 1.
\begin{frame}[plain,label=w2_appA_forms]
  \Large
  \[
    \tag{24}\;\text{Direct form: } Q^S(P), \quad Q^D(P)
  \]
  \[
    \tag{25}\;\text{Inverse form: } P^S(Q), \quad P^D(Q)
  \]
  \[
    \tag{26}\;Q^S = P - 50 \;\Leftrightarrow\; P = 50 + Q
  \]
  \[
    \tag{27}\;Q^D = 85 - 0.5P \;\Leftrightarrow\; P = 170 - 2Q
  \]

  \vspace{0.5em}
  \begin{itemize}\Large
    \item Valid for $P \ge 50$ (supply), $P \le 170$ (demand)
    \item Inverse slope $= 1 /$ direct slope
    \item Demand: $-0.5$ and $-2$; supply: $1$ and $1$
  \end{itemize}
\end{frame}

% --- Appendix C [153]: pass-through in elasticities ------------------------
% Hand check: 2.25 + 1.125 = 3.375; 2.25/3.375 = 0.6667 = 2/3
\begin{frame}[plain,label=w2_appC_passthrough]
  \Large
  \begin{align*}
    \tag{28}\;\rho &= \frac{b}{b + s}
      = \frac{1/|\varepsilon_D|}{1/|\varepsilon_D| + 1/\varepsilon_S}\\
    &= \frac{\varepsilon_S}{\varepsilon_S + |\varepsilon_D|}
  \end{align*}
  \[
    \tag{29}\;\rho = \frac{2.25}{2.25 + 1.125} = \frac{2.25}{3.375} = \frac{2}{3}
  \]

  \vspace{0.3em}
  \begin{itemize}\Large
    \item $b = (P/Q)/|\varepsilon_D|$, $s = (P/Q)/\varepsilon_S$; $P/Q$ cancels
    \item The less elastic side receives more of the saving
  \end{itemize}
\end{frame}

% --- Appendix C [157]: general form of the derived demand curve -------------
% Hand check: P_N = A - c - bQ - gQ = (A - c) - (b + g)Q;
% (220 - 50) - (1 + 1)Q = 170 - 2Q
\begin{frame}[plain,label=w2_appC_derived]
  \Large
  \begin{align*}
    \tag{30}\;A - bQ &= P_N + c + gQ\\
    \Rightarrow P_N &= (A - c) - (b + g)Q
  \end{align*}
  \begin{align*}
    \tag{31}\;P_N &= (220 - 50) - (1 + 1)Q\\
                  &= 170 - 2Q
  \end{align*}

  \vspace{0.3em}
  \begin{itemize}\Large
    \item $-b$: battery demand ($b = 1$)
    \item $-g$: rising other marginal cost ($g = 1$)
    \item $A - c = 220 - 50$: WTP for the first unit of nickel
    \item $A = 170$, $b = 2$ in Walkthrough 4 (nickel)
  \end{itemize}
\end{frame}



%========================
% WEEK 3: Government subsidies and learning-by-doing (Tutorial 3).
% Handout structure: Exercise 1 (set-up, Walkthrough 1 with Questions
% 1.1-1.3, Walkthrough 2 with Question 1.4), Exercise 2 (Questions
% 2.1-2.5, Walkthrough consumer subsidy Steps 1-7, Question 2.6),
% Exercise 3 (A: learning curves, Questions 3.1-3.3; B: Walkthroughs
% 3-5 with Questions 3.4-3.8), Appendix A (general incidence).
% Source: seminar3_handout_with_model_responses.docx (2026-10-19).
% Only formulas from the student-facing text (set-up, walkthroughs,
% Exercise 3 text, Appendix A) are frames; model responses stay text in
% both handout versions (Florian, 19 Oct 2026). One frame per run of
% consecutive equation paragraphs; the image replaces that run. Tags
% (1), (2), ... run in handout order. Export (26 pages, see the header):
%   py build.py peib_main.tex --frames <the w3_ labels below>
%   py export_frames.py peib_main_frames.pdf --out png/week3 --size-factor 0.83 --names <labels>
%========================
\section{Week 3}

% --- Exercise 1 set-up: the cost function ----------------------------------
% Hand check: ten identical firms, F = fixed cost
\begin{frame}[plain,label=w3_setup_tc]
  \Large
  \[
    \tag{1}\;TC_i(q_i) = F + 20\,q_i + 5\,q_i^2
  \]
  \vspace{0.3em}
  \begin{itemize}\Large
    \item $i$: one of the ten firms; $q_i$: its output; $F$: fixed cost
  \end{itemize}
\end{frame}

% --- Exercise 1 set-up: direct demand -----------------------------------------
\begin{frame}[plain,label=w3_setup_demand]
  \Large
  \[
    \tag{2}\;Q^D(P) = 100 - P
  \]
  \vspace{0.3em}
  \begin{itemize}\Large
    \item $Q$: industry quantity ($Q = q_1 + \dots + q_{10}$)
  \end{itemize}
\end{frame}

% --- Walkthrough 1: from the cost function to marginal cost ---------------
% Hand check: d/dq (F + 20q + 5q^2) = 20 + 10q
\begin{frame}[plain,label=w3_wt1_mc]
  \Large
  \[
    TC_i(q_i) = F + 20\,q_i + 5\,q_i^2
  \]
  \[
    \tag{3}\;MC_i(q_i) = \frac{dTC_i}{dq_i} = 20 + 10\,q_i
  \]
  \vspace{0.3em}
  \begin{itemize}\Large
    \item $MC$: the cost of one more panel; $F$ drops out
  \end{itemize}
\end{frame}

% --- Walkthrough 1: the supply rule of a price-taker -------------------------
% Hand check: AVC = 20 + 5q, minimum 20 at q = 0
\begin{frame}[plain,label=w3_wt1_rule]
  \Large
  \[
    \tag{4}\;P = MC_i(q_i) = 20 + 10\,q_i, \quad P \ge 20
  \]
  \vspace{0.3em}
  \begin{itemize}\Large
    \item Price-taker: produce where $P = MC_i$
    \item Below $P = 20$ no panel covers its variable cost
  \end{itemize}
\end{frame}

% --- Walkthrough 2: equilibrium ---------------------------------------------
% Hand check: 100 - Q = 20 + Q => 2Q = 80 => Q = 40; P = 60; 100 - 40 = 60;
% 40/10 = 4 per firm
\begin{frame}[plain,label=w3_wt2_equilibrium]
  \Large
  \begin{align*}
    \tag{5}\;100 - Q &= 20 + Q\\
    \Rightarrow 2Q &= 80\\
    \Rightarrow Q_0 &= 40
  \end{align*}
  \begin{align*}
    \tag{6}\;P_0 &= 20 + 40 = 60\\
    \text{Check: } 100 - 40 &= 60 \;\checkmark
  \end{align*}
  \vspace{0.2em}
  \begin{itemize}\Large
    \item $E_0 = (40, 60)$; per firm $q_i = 40/10 = 4$
  \end{itemize}
\end{frame}

% --- Walkthrough 2: the consumer surplus formula (Question 1.4 text) ------
\begin{frame}[plain,label=w3_wt2_cs_def]
  \Large
  \[
    \tag{7}\;CS = \tfrac{1}{2} \times Q_0 \times \bigl(\text{max.\ willingness to pay} - P_0\bigr)
  \]
  \vspace{0.3em}
  \begin{itemize}\Large
    \item Max.\ willingness to pay $= 100$ (demand intercept)
    \item Added up over all $Q_0$ panels bought
  \end{itemize}
\end{frame}

% --- Walkthrough 2: the producer surplus formula ---------------------------
\begin{frame}[plain,label=w3_wt2_ps_def]
  \Large
  \[
    \tag{8}\;PS = \tfrac{1}{2} \times Q_0 \times \bigl(P_0 - \text{min.\ supply price}\bigr)
  \]
  \vspace{0.3em}
  \begin{itemize}\Large
    \item Min.\ supply price $= 20$ (supply intercept, $MC$ of the first panel)
    \item Added up over all $Q_0$ panels sold
  \end{itemize}
\end{frame}

% --- Walkthrough 2: welfare at E0 -------------------------------------------
% Hand check: 0.5*40*40 = 800 each; 1,600 in total
\begin{frame}[plain,label=w3_wt2_welfare]
  \Large
  \[
    \tag{9}\;CS_0 = \tfrac{1}{2} \times 40 \times (100 - 60) = 800
  \]
  \[
    \tag{10}\;PS_0 = \tfrac{1}{2} \times 40 \times (60 - 20) = 800
  \]
  \vspace{0.3em}
  \begin{itemize}\Large
    \item The two triangles of Figure 1, base $Q_0 = 40$
  \end{itemize}
\end{frame}

% --- Walkthrough 2: total surplus -------------------------------------------
\begin{frame}[plain,label=w3_wt2_ts]
  \Large
  \[
    \tag{11}\;TS_0 = CS_0 + PS_0 = 800 + 800 = 1{,}600
  \]
\end{frame}

% --- Consumer subsidy walkthrough, Step 1: two prices ----------------------
\begin{frame}[plain,label=w3_cs_step1]
  \Large
  \[
    \tag{12}\;P_C = P_P - s = P_P - 20
  \]
  \vspace{0.3em}
  \begin{itemize}\Large
    \item $P_P$: the shop price, received by firms
    \item $P_C$: what consumers pay net of the $20$ they get back
  \end{itemize}
\end{frame}

% --- Step 2: demand in terms of the producer price -------------------------
% Hand check: 100 - (P_P - 20) = 120 - P_P
\begin{frame}[plain,label=w3_cs_step2]
  \Large
  \[
    \tag{13}\;Q^D(P_P) = 100 - (P_P - 20) = 120 - P_P
  \]
  \vspace{0.3em}
  \begin{itemize}\Large
    \item Replace $P_C$ by $P_P - 20$ in $Q^D = 100 - P_C$
    \item At every shop price consumers buy $20$ more panels
  \end{itemize}
\end{frame}

% --- Step 3: inverse demand -------------------------------------------------
\begin{frame}[plain,label=w3_cs_step3]
  \Large
  \[
    \tag{14}\;P_P^D(Q) = 120 - Q
  \]
  \vspace{0.3em}
  \begin{itemize}\Large
    \item The shop price at which consumers buy $Q$ panels
    \item Demand shifts up by $s = 20$ in the $(Q, P_P)$ diagram
  \end{itemize}
\end{frame}

% --- Step 4: supply is unchanged --------------------------------------------
\begin{frame}[plain,label=w3_cs_step4]
  \Large
  \[
    \tag{15}\;P_P^S(Q) = 20 + Q
  \]
  \vspace{0.3em}
  \begin{itemize}\Large
    \item Firms receive $P_P$; costs have not changed (Walkthrough 1)
  \end{itemize}
\end{frame}

% --- Step 5: equilibrium quantity -------------------------------------------
% Hand check: 120 - Q = 20 + Q => 2Q = 100 => Q = 50
\begin{frame}[plain,label=w3_cs_step5]
  \Large
  \begin{align*}
    \tag{16}\;120 - Q &= 20 + Q\\
    \Rightarrow 2Q &= 100\\
    \Rightarrow Q &= 50
  \end{align*}
  \vspace{0.2em}
  \begin{itemize}\Large
    \item Both prices in producer-price terms, so they can be set equal
  \end{itemize}
\end{frame}

% --- Step 6: the two prices ---------------------------------------------------
% Hand check: 20 + 50 = 70; 70 - 20 = 50; demand 100 - 50 = 50
\begin{frame}[plain,label=w3_cs_step6]
  \Large
  \[
    \tag{17}\;P_P = 20 + 50 = 70
  \]
  \begin{align*}
    \tag{18}\;P_C &= 70 - 20 = 50\\
    \text{Check: } 100 - 50 &= 50 \;\checkmark
  \end{align*}
  \vspace{0.2em}
  \begin{itemize}\Large
    \item Same $Q$, $P_C$ and $P_P$ as the producer subsidy (Question 2.3)
  \end{itemize}
\end{frame}

% --- Exercise 3A: the learning curve ----------------------------------------
% Hand check: doubling E multiplies c by 2^(-beta); beta = 0.32: 0.80;
% beta = 1: 0.5
\begin{frame}[plain,label=w3_learning_curve]
  \Large
  \[
    \tag{19}\;c(E) = c_0 \times E^{-\beta}
  \]
  \vspace{0.3em}
  \begin{itemize}\Large
    \item $c_0$: unit cost at the reference experience ($E = 1$)
    \item $E$: cumulative experience relative to that reference
    \item $\beta > 0$: learning parameter; each doubling of $E$ multiplies cost by $2^{-\beta}$
  \end{itemize}
\end{frame}

% --- Exercise 3A: production function with experience ----------------------
\begin{frame}[plain,label=w3_production]
  \Large
  \[
    \tag{20}\;Y_{it} = A_i \times K_{it}^{\alpha} \times L_{it}^{1-\alpha} \times E_{it}^{\gamma}
  \]
  \vspace{0.3em}
  \begin{itemize}\Large
    \item $Y$: output; $K$: capital; $L$: labour; $E$: experience
    \item $A_i$: the firm's productivity level (Lecture 2)
  \end{itemize}
\end{frame}

% --- Exercise 3A: the log form -------------------------------------------------
% Hand check: ln(A K^a L^(1-a) E^g) = ln A + a ln K + (1-a) ln L + g ln E
\begin{frame}[plain,label=w3_production_log]
  \Large
  \begin{align*}
    \tag{21}\;\ln Y_{it} &= \ln A_i + \alpha \ln K_{it} + (1-\alpha) \ln L_{it}\\
                         &\quad + \gamma \ln E_{it}
  \end{align*}
  \vspace{0.2em}
  \begin{itemize}\Large
    \item Products become sums: the form estimated with data
    \item $\gamma$: \% rise in output per 1\% more experience
  \end{itemize}
\end{frame}

% --- Exercise 3A: internal and external learning ----------------------------
\begin{frame}[plain,label=w3_production_spill]
  \Large
  \begin{align*}
    \tag{22}\;\ln Y_{it} &= \ln A_i + \alpha \ln K_{it} + (1-\alpha) \ln L_{it}\\
                         &\quad + \gamma_{int} \ln E_{it} + \gamma_{ext} \ln E^{I}_{t}
  \end{align*}
  \vspace{0.2em}
  \begin{itemize}\Large
    \item $E_{it}$: the firm's own experience; $E^{I}_{t}$: the rest of the industry's
    \item $\gamma_{int}$: internal learning; $\gamma_{ext}$: the spillover ($= 0$ if internal only)
  \end{itemize}
\end{frame}

% --- Exercise 3B set-up: costs at a given level of experience --------------
\begin{frame}[plain,label=w3_cost_E]
  \Large
  \[
    \tag{23}\;TC_i(q_i \mid E) = F + c(E) \times q_i + 5\,q_i^2
  \]
  \[
    \tag{24}\;MC_i(q_i \mid E) = c(E) + 10\,q_i
  \]
  \vspace{0.3em}
  \begin{itemize}\Large
    \item $c(E)$ replaces the intercept $20$; the slope $10$ is unchanged
    \item At $E = 1$: the costs of Exercises 1 and 2
  \end{itemize}
\end{frame}

% --- Walkthrough 3: experience with the subsidy -----------------------------
% Hand check: 50/40 = 1.25
\begin{frame}[plain,label=w3_wt3_experience]
  \Large
  \[
    \tag{25}\;E = \frac{50}{40} = 1.25
  \]
  \vspace{0.3em}
  \begin{itemize}\Large
    \item Period-1 output $50$ with the subsidy, $40$ without
    \item $E = 1$: the no-subsidy path; only the extra output counts
  \end{itemize}
\end{frame}

% --- Walkthrough 3: the new cost intercept ----------------------------------
% Hand check: 20 * 1.25^(-1) = 20/1.25 = 16
\begin{frame}[plain,label=w3_wt3_cost]
  \Large
  \[
    \tag{26}\;c(E) = 20 \times 1.25^{-1} = \frac{20}{1.25} = 16
  \]
  \vspace{0.3em}
  \begin{itemize}\Large
    \item $\beta = 1$: a doubling of $E$ would halve the cost
    \item Every firm's cost intercept falls by $4$
  \end{itemize}
\end{frame}

% --- Walkthrough 4: industry supply in period 2 (no subsidy) ----------------
% Hand check: 10*(P - 16)/10 = P - 16; P = 16 + Q
\begin{frame}[plain,label=w3_wt4_supply]
  \Large
  \[
    \tag{27}\;Q^S(P) = 10 \times \frac{P - 16}{10} = P - 16
  \]
  \[
    \tag{28}\;P^S(Q) = 16 + Q
  \]
  \vspace{0.3em}
  \begin{itemize}\Large
    \item Exercise 1 had $P = 20 + Q$: the curve shifts down by the saving of $4$
  \end{itemize}
\end{frame}

% --- Walkthrough 5: surpluses and the ledger in period 2 (no subsidy) ------
% Hand check: E = (42, 58); 0.5*42*42 = 882 each; G = 0; 1,764
\begin{frame}[plain,label=w3_wt5_ledger]
  \Large
  \[
    \tag{29}\;CS = \tfrac{1}{2} \times 42 \times (100 - 58) = 882
  \]
  \[
    \tag{30}\;PS = \tfrac{1}{2} \times 42 \times (58 - 16) = 882
  \]
  \[
    \tag{31}\;G = 0
  \]
  \begin{align*}
    \tag{32}\;CS + PS - G &= 882 + 882 - 0\\
                 &= 1{,}764
  \end{align*}
\end{frame}

% --- Appendix A: who gets the subsidy, general linear case -----------------
% Hand check: a - bQ = c + dQ - s => Q1 = (a - c + s)/(b + d);
% dQ = s/(b + d); dP_C = -b dQ; dP_P = dP_C + s = d s/(b + d)
\begin{frame}[plain,label=w3_appA_incidence]
  \Large
  \begin{align*}
    \tag{33}\;\Delta Q &= \frac{s}{b + d}\\
    \Delta P_C &= -\frac{b\,s}{b + d}\\
    \Delta P_P &= +\frac{d\,s}{b + d}
  \end{align*}
  \vspace{0.3em}
  \begin{itemize}\Large
    \item Consumers' share $b/(b + d)$, producers' share $d/(b + d)$
    \item The steeper (less elastic) side receives more
  \end{itemize}
\end{frame}

% --- Appendix A: the handout numbers ------------------------------------------
% Hand check: b = d = 1, s = 20: dQ = 10, dP_C = -10, dP_P = +10
\begin{frame}[plain,label=w3_appA_example]
  \Large
  \[
    \tag{34}\;\Delta Q = 20/2 = 10, \quad \Delta P_C = -10, \quad \Delta P_P = +10
  \]
  \vspace{0.3em}
  \begin{itemize}\Large
    \item $b = d = 1$ and $s = 20$, as in Question 2.3
    \item With $s < 0$ the same formulas describe a per-unit tax
  \end{itemize}
\end{frame}

\end{document}
